% \chapter{CƠ SỞ LÝ THUYẾT VÀ CÁC NGHIÊN CỨU LIÊN QUAN}

% \section{Tổng quan về tín dụng và rủi ro tín dụng}

% \subsection{Khái niệm tín dụng}

% Tín dụng là quan hệ kinh tế giữa bên cho vay và bên đi vay, trong đó bên cho vay chuyển giao quyền sử dụng một lượng giá trị nhất định cho bên đi vay trong một khoảng thời gian xác định, với cam kết hoàn trả cả gốc và lãi theo thỏa thuận. Trong hệ thống ngân hàng thương mại, tín dụng là hoạt động kinh doanh cốt lõi, tạo ra phần lớn doanh thu và lợi nhuận cho ngân hàng \cite{thomas2000survey}.

% Các hình thức tín dụng phổ biến bao gồm:

% \begin{itemize}
%     \item Cho vay ngắn hạn;
%     \item Cho vay trung và dài hạn;
%     \item Cho vay tiêu dùng;
%     \item Cho vay sản xuất kinh doanh;
%     \item Thẻ tín dụng;
%     \item Bảo lãnh ngân hàng.
% \end{itemize}

% Sự phát triển của hoạt động tín dụng góp phần thúc đẩy tăng trưởng kinh tế, mở rộng sản xuất kinh doanh và nâng cao mức sống của người dân. Tuy nhiên, đi kèm với lợi nhuận là các rủi ro tiềm ẩn, trong đó rủi ro tín dụng được xem là loại rủi ro quan trọng nhất đối với ngân hàng thương mại \cite{basel2000creditrisk}.

% \subsection{Khái niệm rủi ro tín dụng}

% Theo Basel Committee on Banking Supervision, rủi ro tín dụng là khả năng tổn thất phát sinh do khách hàng hoặc đối tác không thực hiện đầy đủ nghĩa vụ tài chính theo các điều khoản đã cam kết \cite{basel2000creditrisk}.

% Rủi ro tín dụng có thể phát sinh từ nhiều nguyên nhân khác nhau như:

% \begin{itemize}
%     \item Khả năng tài chính yếu kém của khách hàng;
%     \item Suy giảm hoạt động kinh doanh;
%     \item Biến động kinh tế vĩ mô;
%     \item Thiếu thông tin trong quá trình thẩm định;
%     \item Gian lận hoặc cố ý chiếm dụng vốn.
% \end{itemize}

% Hậu quả của rủi ro tín dụng không chỉ ảnh hưởng đến khả năng sinh lời của ngân hàng mà còn có thể đe dọa sự ổn định của toàn bộ hệ thống tài chính \cite{basel_framework}.

% \subsection{Phân loại nợ trong ngân hàng thương mại}

% Theo quy định của Ngân hàng Nhà nước Việt Nam, các khoản nợ được phân thành năm nhóm tương ứng với các mức độ rủi ro khác nhau \cite{sbv2021tt11}.

% \begin{table}[H]
% \centering
% \caption{Phân loại nhóm nợ theo quy định của Ngân hàng Nhà nước}
% \begin{tabular}{|c|p{8cm}|}
% \hline
% Nhóm nợ & Mô tả \\
% \hline
% Nhóm 1 & Nợ đủ tiêu chuẩn \\
% \hline
% Nhóm 2 & Nợ cần chú ý \\
% \hline
% Nhóm 3 & Nợ dưới tiêu chuẩn \\
% \hline
% Nhóm 4 & Nợ nghi ngờ \\
% \hline
% Nhóm 5 & Nợ có khả năng mất vốn \\
% \hline
% \end{tabular}
% \end{table}

% Việc phân loại nợ là cơ sở để ngân hàng thực hiện trích lập dự phòng rủi ro, đánh giá chất lượng tín dụng và xây dựng các mô hình cảnh báo sớm.

% \section{Xếp hạng tín dụng}

% \subsection{Khái niệm xếp hạng tín dụng}

% Xếp hạng tín dụng (Credit Rating hoặc Credit Scoring) là quá trình đánh giá mức độ tín nhiệm và khả năng thực hiện nghĩa vụ tài chính của khách hàng dựa trên các thông tin tài chính và phi tài chính \cite{hand1997statistical,thomas2000survey}.

% Mục tiêu của xếp hạng tín dụng bao gồm:

% \begin{itemize}
%     \item Hỗ trợ quyết định cấp tín dụng;
%     \item Xác định hạn mức tín dụng;
%     \item Xác định lãi suất cho vay;
%     \item Cảnh báo sớm rủi ro tín dụng;
%     \item Hỗ trợ phân loại nợ.
% \end{itemize}

% \subsection{Vai trò của xếp hạng tín dụng}

% Xếp hạng tín dụng đóng vai trò quan trọng trong quản trị rủi ro ngân hàng:

% \begin{enumerate}
%     \item Giảm thiểu rủi ro tín dụng.
%     \item Chuẩn hóa quy trình ra quyết định.
%     \item Tối ưu hóa danh mục tín dụng.
%     \item Hỗ trợ tuân thủ Basel II và Basel III.
% \end{enumerate}

% Theo Basel II, hệ thống xếp hạng nội bộ là một thành phần quan trọng trong phương pháp tiếp cận IRB (Internal Rating Based Approach) để tính toán yêu cầu vốn đối với rủi ro tín dụng \cite{basel_framework}.

% \subsection{Các phương pháp xếp hạng tín dụng truyền thống}

% \subsubsection{Phương pháp chuyên gia}

% Phương pháp này dựa trên kinh nghiệm và đánh giá của cán bộ tín dụng. Ưu điểm là dễ triển khai nhưng mang tính chủ quan cao.

% \subsubsection{Mô hình điểm số tín dụng}

% Một trong những mô hình nổi tiếng là Z-score của Altman \cite{altman1968financial}.

% Mô hình có dạng:

% \[
% Z = 1.2X_1 + 1.4X_2 + 3.3X_3 + 0.6X_4 + 1.0X_5
% \]

% trong đó:

% \begin{itemize}
%     \item $X_1$: Vốn lưu động/Tổng tài sản.
%     \item $X_2$: Lợi nhuận giữ lại/Tổng tài sản.
%     \item $X_3$: EBIT/Tổng tài sản.
%     \item $X_4$: Giá trị thị trường vốn chủ sở hữu/Tổng nợ.
%     \item $X_5$: Doanh thu/Tổng tài sản.
% \end{itemize}

% Mặc dù có giá trị thực tiễn cao, mô hình này gặp hạn chế khi áp dụng cho dữ liệu lớn và dữ liệu phi tuyến.

% \section{Trí tuệ nhân tạo và học máy trong xếp hạng tín dụng}

% \subsection{Tổng quan về trí tuệ nhân tạo}

% Trí tuệ nhân tạo (Artificial Intelligence - AI) là lĩnh vực nghiên cứu nhằm phát triển các hệ thống có khả năng thực hiện các nhiệm vụ đòi hỏi trí tuệ của con người như học tập, suy luận, nhận dạng và ra quyết định \cite{worldbank2025ai}.

% Trong ngành tài chính, AI được ứng dụng trong:

% \begin{itemize}
%     \item Chấm điểm tín dụng;
%     \item Phát hiện gian lận;
%     \item Chống rửa tiền;
%     \item Dự báo rủi ro;
%     \item Cá nhân hóa dịch vụ.
% \end{itemize}

% \subsection{Machine Learning}

% Machine Learning là một nhánh của AI cho phép máy tính học từ dữ liệu mà không cần lập trình tường minh \cite{bishop2006pattern}.

% Machine Learning được chia thành:

% \begin{itemize}
%     \item Học có giám sát (Supervised Learning);
%     \item Học không giám sát (Unsupervised Learning);
%     \item Học tăng cường (Reinforcement Learning).
% \end{itemize}

% Trong nghiên cứu này, bài toán xếp hạng tín dụng thuộc nhóm học có giám sát.

% \section{Các thuật toán học máy sử dụng trong nghiên cứu}

% \subsection{Logistic Regression}

% Logistic Regression là một trong những thuật toán phổ biến nhất trong bài toán tín dụng \cite{hosmer2013applied}.

% Xác suất xảy ra sự kiện được xác định bởi:

% \[
% P(Y=1|X)
% =
% \frac{1}
% {1+e^{-(\beta_0+\beta_1x_1+\cdots+\beta_nx_n)}}
% \]

% Ưu điểm:

% \begin{itemize}
%     \item Dễ triển khai.
%     \item Dễ giải thích.
%     \item Hiệu quả với dữ liệu tuyến tính.
% \end{itemize}

% Nhược điểm:

% \begin{itemize}
%     \item Khó mô hình hóa quan hệ phi tuyến.
% \end{itemize}

% \subsection{Decision Tree}

% Decision Tree xây dựng tập luật phân chia dữ liệu theo cấu trúc cây \cite{quinlan1986induction}.

% Ưu điểm:

% \begin{itemize}
%     \item Dễ giải thích.
%     \item Không yêu cầu chuẩn hóa dữ liệu.
% \end{itemize}

% Nhược điểm:

% \begin{itemize}
%     \item Dễ overfitting.
% \end{itemize}

% \subsection{Random Forest}

% Random Forest là phương pháp ensemble được đề xuất bởi Breiman \cite{breiman2001random}.

% Ý tưởng chính là xây dựng nhiều cây quyết định trên các tập dữ liệu bootstrap khác nhau.

% Ưu điểm:

% \begin{itemize}
%     \item Giảm overfitting.
%     \item Độ chính xác cao.
% \end{itemize}

% \subsection{XGBoost}

% XGBoost là thuật toán boosting được sử dụng rộng rãi trong các bài toán dự báo tín dụng \cite{chen2016xgboost}.

% Hàm mục tiêu:

% \[
% Obj =
% \sum_i l(y_i,\hat{y}_i)
% +
% \sum_k \Omega(f_k)
% \]

% Trong đó:

% \[
% \Omega(f)
% =
% \gamma T
% +
% \frac{1}{2}
% \lambda ||w||^2
% \]

% Thuật toán có khả năng xử lý dữ liệu lớn và cho hiệu quả dự báo rất cao.

% \subsection{LightGBM}

% LightGBM được Microsoft phát triển nhằm tối ưu tốc độ huấn luyện và khả năng mở rộng của Gradient Boosting \cite{ke2017lightgbm}.

% Ưu điểm:

% \begin{itemize}
%     \item Tốc độ nhanh.
%     \item Bộ nhớ thấp.
%     \item Hiệu quả trên dữ liệu lớn.
% \end{itemize}

% \subsection{CatBoost}

% CatBoost được phát triển bởi Yandex nhằm xử lý tốt dữ liệu chứa biến phân loại \cite{prokhorenkova2018catboost}.

% Ưu điểm:

% \begin{itemize}
%     \item Xử lý categorical feature hiệu quả.
%     \item Giảm overfitting.
%     \item Độ chính xác cao.
% \end{itemize}

% \section{Các chỉ tiêu đánh giá mô hình}

% Các mô hình được đánh giá bằng các chỉ tiêu phổ biến trong bài toán phân loại \cite{fawcett2006roc}.

% \subsection{Accuracy}

% \[
% Accuracy
% =
% \frac{TP+TN}
% {TP+TN+FP+FN}
% \]

% \subsection{Precision}

% \[
% Precision
% =
% \frac{TP}
% {TP+FP}
% \]

% \subsection{Recall}

% \[
% Recall
% =
% \frac{TP}
% {TP+FN}
% \]

% \subsection{F1-Score}

% \[
% F1
% =
% 2\times
% \frac{Precision \times Recall}
% {Precision + Recall}
% \]

% \section{Tổng quan các nghiên cứu liên quan}

% \subsection{Nghiên cứu quốc tế}

% Các nghiên cứu của Hand và Henley \cite{hand1997statistical}, Thomas \cite{thomas2000survey}, Baesens \cite{baesens2003benchmarking} và Lessmann \cite{lessmann2015benchmarking} cho thấy các phương pháp học máy thường cho hiệu quả vượt trội hơn các phương pháp thống kê truyền thống trong bài toán chấm điểm tín dụng.

% Gần đây, các thuật toán boosting như XGBoost và LightGBM đã đạt được nhiều kết quả khả quan trong các cuộc thi và ứng dụng thực tế liên quan đến dự báo rủi ro tín dụng \cite{chen2016xgboost,ke2017lightgbm}.

% \subsection{Nghiên cứu trong nước}

% Tại Việt Nam, các nghiên cứu về xếp hạng tín dụng chủ yếu tập trung vào:

% \begin{itemize}
%     \item Chấm điểm tín dụng cá nhân.
%     \item Đánh giá rủi ro doanh nghiệp.
%     \item Dự báo nợ xấu.
% \end{itemize}

% Tuy nhiên, số lượng nghiên cứu sử dụng dữ liệu thực tế của ngân hàng thương mại và áp dụng đồng thời nhiều thuật toán học máy hiện đại vẫn còn hạn chế.

% \section{Khoảng trống nghiên cứu}

% Từ tổng quan tài liệu có thể nhận thấy:

% \begin{itemize}
%     \item Thiếu nghiên cứu sử dụng dữ liệu tín dụng thực tế tại Việt Nam.
%     \item Thiếu nghiên cứu so sánh đồng thời nhiều thuật toán AI hiện đại.
%     \item Chưa có nhiều nghiên cứu đánh giá khả năng ứng dụng thực tiễn tại ngân hàng thương mại Việt Nam.
% \end{itemize}

% Do đó, đề tài này được thực hiện nhằm bổ sung khoảng trống nghiên cứu trên thông qua việc ứng dụng các mô hình trí tuệ nhân tạo trong xếp hạng tín dụng dựa trên dữ liệu thực tế của một ngân hàng thương mại tại Việt Nam.



\chapter{CƠ SỞ LÝ THUYẾT VÀ CÁC NGHIÊN CỨU LIÊN QUAN}

\section{Tín dụng và rủi ro tín dụng}

\subsection{Khái niệm tín dụng}

Tín dụng là quan hệ kinh tế trong đó bên cho vay chuyển giao quyền sử dụng một lượng giá trị cho bên đi vay trong một khoảng thời gian xác định, với nghĩa vụ hoàn trả theo các điều kiện đã thỏa thuận.

\subsection{Khái niệm rủi ro tín dụng}

Rủi ro tín dụng là khả năng phát sinh tổn thất do khách hàng hoặc đối tác không thực hiện đầy đủ nghĩa vụ tài chính. Rủi ro chịu ảnh hưởng từ năng lực tài chính, đặc điểm khoản vay, lịch sử tín dụng, chính sách cấp tín dụng và điều kiện kinh tế.

\subsection{Phân loại nhóm nợ}

\begin{table}[H]
\centering
\caption{Các nhóm nợ trong nghiên cứu}
\begin{tabular}{|c|p{4cm}|p{7cm}|}
\hline
1 & Nợ đủ tiêu chuẩn & Mức rủi ro thấp.\\
\hline
2 & Nợ cần chú ý & Đã xuất hiện dấu hiệu cần theo dõi.\\
\hline
3 & Nợ dưới tiêu chuẩn & Khả năng thu hồi suy giảm đáng kể.\\
\hline
4 & Nợ nghi ngờ & Mức rủi ro cao.\\
\hline
5 & Nợ có khả năng mất vốn & Mức rủi ro cao nhất.\\
\hline
\end{tabular}
\end{table}

\section{Chấm điểm và xếp hạng tín dụng}

Chấm điểm tín dụng tạo ra một giá trị định lượng phản ánh rủi ro; xếp hạng tín dụng gán khách hàng hoặc khoản vay vào một nhóm rủi ro. Hai khái niệm liên hệ nhưng không đồng nhất \cite{hand1997statistical,thomas2000survey}.

Tổn thất kỳ vọng có thể biểu diễn:
\[
EL=PD\times LGD\times EAD.
\]

\section{Các phương pháp đánh giá tín dụng truyền thống}

Phương pháp chuyên gia có khả năng khai thác thông tin định tính nhưng chủ quan. Z-score của Altman có ý nghĩa nền tảng nhưng phụ thuộc dữ liệu kế toán và quan hệ tuyến tính. Logistic Regression dễ giải thích nhưng khó tự động nắm bắt tương tác phức tạp.

\section{Trí tuệ nhân tạo và học máy}

Học máy là nhánh của AI trong đó mô hình học quy luật từ dữ liệu. Bài toán của luận văn thuộc học có giám sát.

\section{Các mô hình sử dụng}

\subsection{Logistic Regression đa lớp}

\[
P(Y=k\mid\mathbf{x})
=
\frac{\exp(\beta_{k0}+\boldsymbol{\beta}_{k}^{T}\mathbf{x})}
{\sum_{j=1}^{K}\exp(\beta_{j0}+\boldsymbol{\beta}_{j}^{T}\mathbf{x})}.
\]

\subsection{Decision Tree}

Decision Tree chia không gian đặc trưng bằng các điều kiện dạng $x_j\le s$. Mô hình dễ giải thích nhưng dễ overfitting.

\subsection{Random Forest}

\[
\hat{y}=\operatorname{mode}\{h_b(\mathbf{x})\}_{b=1}^{B}.
\]

\subsection{XGBoost}

\[
\mathcal{L}^{(t)}
=
\sum_{i=1}^{n}
l(y_i,\hat{y}_{i}^{(t-1)}+f_t(\mathbf{x}_i))
+\Omega(f_t).
\]

\subsection{LightGBM}

LightGBM tối ưu tốc độ và bộ nhớ, nhưng cần kiểm soát độ sâu và số lá.

\subsection{CatBoost}

CatBoost xử lý biến phân loại và sử dụng ordered boosting để hạn chế leakage.

\section{Các vấn đề đặc thù của dữ liệu tín dụng}

\subsection{Mất cân bằng lớp}

\[
\mathcal{L}
=
-\sum_{i=1}^{N}
w_{y_i}
\sum_{k=1}^{K}
\mathbb{I}(y_i=k)\log p_{ik}.
\]

\subsection{Tính thứ tự}

\[
MAE_{\mathrm{ordinal}}
=
\frac{1}{N}\sum_{i=1}^{N}|y_i-\hat{y}_i|.
\]

\subsection{Thay đổi theo thời gian}

\[
\mathcal{D}_{train}<\mathcal{D}_{validation}<\mathcal{D}_{test}.
\]

\subsection{Khả năng giải thích}

\[
f(\mathbf{x})=\phi_0+\sum_{j=1}^{p}\phi_j.
\]

\section{Chỉ tiêu đánh giá}

\subsection{Accuracy}
\[
Accuracy=
\frac{\sum_{k=1}^{K}C_{kk}}
{\sum_{i=1}^{K}\sum_{j=1}^{K}C_{ij}}.
\]

\subsection{Precision, Recall và F1}
\[
Precision_k=\frac{TP_k}{TP_k+FP_k},
\qquad
Recall_k=\frac{TP_k}{TP_k+FN_k},
\]
\[
F1_k=2\frac{Precision_k\times Recall_k}{Precision_k+Recall_k}.
\]

\subsection{Macro F1}
\[
F1_{\mathrm{macro}}=\frac{1}{K}\sum_{k=1}^{K}F1_k.
\]

\subsection{Weighted F1}
\[
F1_{\mathrm{weighted}}=\sum_{k=1}^{K}\frac{n_k}{N}F1_k.
\]

\subsection{Balanced Accuracy}
\[
Balanced\ Accuracy=\frac{1}{K}\sum_{k=1}^{K}Recall_k.
\]

\subsection{Ordinal MAE}
\[
Ordinal\ MAE=\frac{1}{N}\sum_{i=1}^{N}|y_i-\hat{y}_i|.
\]

\section{Tổng quan nghiên cứu liên quan}

Hand và Henley \cite{hand1997statistical} hệ thống hóa các phương pháp thống kê; Thomas \cite{thomas2000survey} mở rộng từ application scoring sang behavioural scoring. Baesens và cộng sự \cite{baesens2003benchmarking} và Lessmann và cộng sự \cite{lessmann2015benchmarking} cho thấy hiệu quả mô hình phụ thuộc vào dữ liệu, thước đo và quy trình đánh giá.

Brown và Mues \cite{brown2012experimental} nhấn mạnh ảnh hưởng của mất cân bằng. Xia và cộng sự \cite{xia2017boosted} cho thấy boosted tree có khả năng cải thiện dự báo. Barboza và cộng sự \cite{barboza2017machine} ghi nhận kết quả cạnh tranh của các mô hình cây trong dự báo phá sản. Dastile và cộng sự \cite{dastile2020statistical} chỉ ra các hạn chế phổ biến về dữ liệu nhỏ, mất cân bằng, thiếu chuẩn đánh giá và khả năng giải thích.

Tại Việt Nam, nghiên cứu thường tập trung vào credit scoring cá nhân, rủi ro doanh nghiệp và dự báo nợ xấu. Tuy nhiên, chưa có nhiều nghiên cứu công bố việc phân loại trực tiếp năm nhóm nợ trên dữ liệu thực tế, kết hợp so sánh nhiều mô hình, temporal split, ordinal evaluation và explainability.

\section{Khoảng trống nghiên cứu}

\begin{enumerate}
    \item Phần lớn nghiên cứu dùng mục tiêu nhị phân.
    \item Tính thứ tự của năm nhóm nợ chưa được khai thác đầy đủ.
    \item Dữ liệu thực tế tại ngân hàng Việt Nam ít được công bố.
    \item Random split vẫn phổ biến hơn temporal split.
    \item Lớp rủi ro cao ít quan sát nhưng chưa được đánh giá đầy đủ.
    \item Mô hình phức tạp thường thiếu khả năng giải thích.
\end{enumerate}

\section{Định hướng nghiên cứu}

Luận văn xây dựng bài toán năm nhóm nợ, kiểm tra và xây dựng đặc trưng nghiệp vụ, sử dụng Logistic Regression làm baseline, so sánh với mô hình cây và boosting, xử lý mất cân bằng, chia theo thời gian, đánh giá bằng chỉ tiêu đa lớp và thứ tự, đồng thời giải thích mô hình bằng feature importance và SHAP.

\section{Tóm tắt chương}

Chương này đã trình bày cơ sở lý thuyết, thành tựu và hạn chế của các công trình nền tảng, các mô hình được sử dụng, các vấn đề đặc thù của dữ liệu tín dụng, tổng quan nghiên cứu, khoảng trống và định hướng của luận văn.
